% ============================================================
\chapter[Sistemas de solución de conflictos]{Sistemas de solución de conflictos}
% ============================================================

\ideacentral{El proceso judicial es un molde único; por eso existen los métodos alternativos de resolución de conflictos (MASC), que complementan —no sustituyen— a la jurisdicción. Se ordenan según si el tercero tiene o no poder para imponer una solución.}

\section{El conflicto y el mapa de sistemas}

Siguiendo a Remo Entelman, el conflicto es la interferencia intersubjetiva de intereses entre dos o más partes (por ejemplo, el incumplimiento de una locación). Conviene trabajarlo temprano, antes de que escale, tratando el vínculo como un canal de comunicación.

El derecho procesal es \textbf{instrumental}: sirve al derecho de fondo. Las formas dan certeza, pero son una garantía y no un fin en sí mismas.

\begin{center}
\begin{tabularx}{\textwidth}{>{\raggedright\arraybackslash}p{3.4cm} >{\raggedright\arraybackslash}p{2.6cm} >{\raggedright\arraybackslash}p{2.4cm} Y}
\toprule
\textbf{Sistema} & \textbf{Tercero} & \textbf{¿Impone solución?} & \textbf{Ejemplos} \\
\midrule
\textbf{Autocomposición (no adversarial)} & Ninguno & No & Negociación entre partes o abogados \\
\addlinespace
\textbf{Heterocomposición asistida (no adversarial)} & Mediador / conciliador & No & Mediación, conciliación \\
\addlinespace
\textbf{Adversarial} & Árbitro / juez & \textbf{Sí} & Arbitraje, proceso judicial \\
\bottomrule
\end{tabularx}
\end{center}

\section{La negociación por intereses}

\begin{itemize}
  \item \textbf{Separar a las personas del problema.} Los abogados no pueden pelearse con el colega en cada caso.
  \item \textbf{Negociar sobre intereses, no sobre posiciones.} Detrás de «esto vale 50» hay un interés a descubrir.
  \item \textbf{Generar opciones de beneficio mutuo} antes de cerrar una posición (tormenta de ideas).
  \item \textbf{Usar criterios objetivos:} valor de mercado, antecedentes, tasación.
\end{itemize}

\begin{example}{La naranja}
Dos empresas se disputaban el 100\% de una naranja; al indagar los intereses, una necesitaba la cáscara y la otra el jugo.
\end{example}

\subsection*{Límite: las vías de hecho}

Aunque no le paguen, el locador \textbf{no puede recuperar la tenencia por la fuerza} (por ejemplo, poniendo un candado). Hacerlo lo expone a denuncia penal por usurpación y a restituir la tenencia. Solo en una ocupación ilegal recién consumada cabe actuar en defensa de la posesión; cuanto más se consolida la ocupación, menos margen hay.

\section{Mediación y conciliación}

\begin{center}
\begin{tabularx}{\textwidth}{>{\raggedright\arraybackslash}p{2.8cm} Y Y}
\toprule
 & \textbf{Mediador} & \textbf{Conciliador} \\
\midrule
\textbf{Rol} & Acerca intereses para que las partes resuelvan & Puede \textbf{proponer fórmulas conciliatorias} \\
\addlinespace
\textbf{Decide el caso} & No & No \\
\addlinespace
\textbf{Ejemplo} & Mediación prejudicial & El juez en la audiencia del art.\ 360 CPCCN (proponer fórmulas no es prejuzgar) \\
\bottomrule
\end{tabularx}
\end{center}

\textbf{Régimen de la mediación prejudicial obligatoria} (según la clase 3):

\begin{itemize}
  \item Puede ser \textbf{privada} (el requirente elige un mediador de la lista del Ministerio de Justicia y propone tres al requerido) u \textbf{oficial} (formulario ante la cámara; el mediador sale por sorteo).
  \item Dura \textbf{60 días}, prorrogables. Rige la \textbf{confidencialidad}: lo dicho no puede usarse como prueba ni se cita al mediador como testigo.
  \item Si no hay acuerdo, se entrega el \textbf{acta de cierre}, requisito de admisibilidad de la demanda.
  \item La multa por incomparecencia pasó de obligatoria a facultativa («podrá») con la Ley 26.589; en la práctica casi no se aplica.
  \item \textbf{Prescripción:} el inicio de la mediación la suspende hasta \textbf{20 días después del acta de cierre}; luego corre el remanente. Si no se demanda a tiempo, prescribe.
  \item El acuerdo de mediación no requiere homologación (salvo menores o incapaces) y se ejecuta con el acta del mediador.
\end{itemize}

% TODO: contenido incompleto o ambiguo en las notas originales
\begin{important}
\textbf{Inconsistencia a verificar.} La clase 3 describe la mediación como \textbf{obligatoria}, pero las clases 2 y 8 afirman que en el \textbf{desalojo es optativa}. El programa menciona «mediación previa» en el Taller 1 con referencia al DNU 70/23. Conviene confirmar en la norma vigente qué se exige para desalojo.
\end{important}

\section{Arbitraje}

\begin{itemize}
  \item \textbf{Cláusula compromisoria:} se pacta \textbf{antes} del conflicto (habitualmente en el contrato). Tiene función \textbf{positiva} (si la otra parte no firma el compromiso, se puede pedir judicialmente que se constituya y el juez suple su voluntad) y \textbf{negativa} (si te demandan en la justicia, oponés la excepción).
  \item \textbf{Compromiso arbitral:} se firma \textbf{cuando ya existe} el conflicto y fija las cuestiones concretas que resolverá el árbitro, dentro del plazo fijado.
  \item \textbf{Arbitraje de derecho} (fundado en normas) y de \textbf{amigable composición} (equidad: «a verdad sabida y buena fe guardada»).
  \item \textbf{Pericia arbitral:} un perito resuelve con carácter vinculante una cuestión técnica.
  \item \textbf{Límite clave:} el árbitro \textbf{carece de imperium}: no ejecuta su laudo ni dicta cautelares. Para el lanzamiento hay que ir al juez que hubiera sido competente.
\end{itemize}

\section{¿Qué camino elegir?}

\begin{center}
\begin{tabularx}{\textwidth}{>{\raggedright\arraybackslash}p{3.4cm} Y Y}
\toprule
\textbf{Camino} & \textbf{Ventajas} & \textbf{Límites} \\
\midrule
\textbf{Proceso judicial} & Decisión ejecutable por el propio juez, incluido el lanzamiento & Mayor costo y tiempo \\
\addlinespace
\textbf{Arbitraje} & Puede acortar el fondo & Paso adicional ante el juez para ejecutar \\
\addlinespace
\textbf{Mediación / negociación} & Suele ser lo más eficiente en tiempo y costo & Depende de la voluntad de ambas partes \\
\bottomrule
\end{tabularx}
\end{center}

% TODO: contenido incompleto o ambiguo en las notas originales
\begin{note}
\textbf{Vacío parcial.} El programa habla de «diversos regímenes» y «arbitraje de derecho y amigable componedor». Los apuntes tratan el arbitraje de forma conceptual; no desarrollan regímenes específicos más allá de la Ley 26.589 y del arbitraje de colegios profesionales mencionado como ejemplo.
\end{note}

\section{Cuadro Cornell}

\begin{longtable}{Q R}
\toprule
\textbf{Preguntas / palabras clave} & \textbf{Notas / respuestas} \\
\midrule
\endfirsthead
\toprule
\textbf{Preguntas / palabras clave} & \textbf{Notas / respuestas} \\
\midrule
\endhead

\textbf{¿Qué criterio separa los distintos sistemas de solución de conflictos?} &
Si el tercero tiene o no poder para imponer una solución. Autocomposición: sin tercero (negociación). Heterocomposición asistida: mediador o conciliador, que no imponen solución. Adversarial: árbitro o juez, que sí la imponen. \\
\addlinespace[0.6em]

\textbf{¿Cuáles son los cuatro principios de la negociación por intereses?} &
Separar a las personas del problema; negociar sobre intereses y no sobre posiciones; generar opciones de beneficio mutuo antes de cerrar una posición; y usar criterios objetivos (valor de mercado, antecedentes, tasación). \\
\addlinespace[0.6em]

\textbf{¿Por qué el locador no puede recuperar la tenencia por la fuerza?} &
Porque se expone a denuncia penal por usurpación y a restituir la tenencia. Solo en una ocupación ilegal recién consumada cabe actuar en defensa de la posesión; cuanto más se consolida la ocupación, menos margen hay. \\
\addlinespace[0.6em]

\textbf{¿En qué se diferencian mediador y conciliador?} &
El mediador acerca intereses para que las partes resuelvan; el conciliador puede proponer fórmulas conciliatorias. Ninguno decide el caso. Ejemplo de conciliador: el juez en la audiencia del art.\ 360 CPCCN. \\
\addlinespace[0.6em]

\textbf{¿Cómo funciona la mediación prejudicial según la clase 3?} &
Puede ser privada u oficial; dura 60 días, prorrogables; rige la confidencialidad; si no hay acuerdo se entrega el acta de cierre, requisito de admisibilidad de la demanda. La multa por incomparecencia pasó a ser facultativa con la Ley 26.589. El acuerdo no requiere homologación (salvo menores o incapaces) y se ejecuta con el acta del mediador. \\
\addlinespace[0.6em]

\textbf{¿Qué efecto tiene la mediación sobre la prescripción?} &
El inicio de la mediación la suspende hasta 20 días después del acta de cierre; luego corre el remanente. Si no se demanda a tiempo, prescribe. \\
\addlinespace[0.6em]

\textbf{¿Qué diferencia hay entre cláusula compromisoria y compromiso arbitral?} &
La cláusula compromisoria se pacta antes del conflicto, habitualmente en el contrato; tiene función positiva (el juez suple la voluntad de quien no firma el compromiso) y negativa (se opone como excepción si hay demanda judicial). El compromiso arbitral se firma cuando el conflicto ya existe y fija las cuestiones que resolverá el árbitro, dentro del plazo fijado. \\
\addlinespace[0.6em]

\textbf{¿Por qué el árbitro necesita al juez?} &
Porque carece de imperium: no ejecuta su laudo ni dicta cautelares. Para el lanzamiento hay que ir al juez que hubiera sido competente. \\
\addlinespace[0.6em]

\textbf{¿Qué ventajas y límites tiene cada camino?} &
Proceso judicial: decisión ejecutable por el propio juez, pero mayor costo y tiempo. Arbitraje: puede acortar el fondo, pero exige un paso adicional ante el juez para ejecutar. Mediación o negociación: suele ser lo más eficiente, pero depende de la voluntad de ambas partes. \\

\midrule
\multicolumn{2}{p{0.92\textwidth}}{%
\textbf{Resumen:} Lo que separa los sistemas es si el tercero \textbf{impone} una solución. Cuatro principios de negociación: personas/problema, intereses/posiciones, opciones, criterios objetivos. Mediador no propone fórmulas; conciliador sí; ninguno decide. Mediación suspende la prescripción hasta 20 días tras el acta de cierre. Cláusula compromisoria (antes) $\neq$ compromiso arbitral (después). El árbitro no tiene imperium.} \\
\bottomrule
\end{longtable}

% ============================================================
\chapter[Diligencias preliminares y proceso]{Diligencias preliminares y proceso de desalojo}
% ============================================================

\ideacentral{El desalojo es un proceso especial que tramita por conocimiento (ordinario o sumarísimo) y se apoya en toda la teoría procesal general. Antes de demandar puede ser necesario preparar el proceso con diligencias preliminares, sobre todo para identificar al ocupante.}

\section{Diligencias preliminares y medida preparatoria}

\begin{itemize}
  \item Son medidas que pide quien pretende demandar (o prevé, con fundamento, que será demandado) para \textbf{preparar el proceso de conocimiento} (art.\ 323 CPCCN y ss.). Son de carácter restrictivo y residual.
  \item \textbf{Art.\ 323, inc.\ 1:} permite pedir que la persona contra quien se dirigirá la demanda preste \textbf{declaración jurada, por escrito, sobre hechos relativos a su personalidad}: herramienta clave para identificar al ocupante.
  \item No se usan para averiguar el domicilio, que se ubica por medios propios.
  \item \textbf{Plazo de caducidad: 30 días} desde su sustanciación para iniciar el proceso principal; si no, pierde efectos.
  \item \textbf{No es prueba anticipada:} esta es un sistema asegurativo para evitar la pérdida de una prueba (por ejemplo, un testigo próximo a fallecer) y es un instituto autónomo.
\end{itemize}

Comparado: en Estados Unidos (discovery) y Gran Bretaña (disclosure) hay una preparación con fuerte componente extrajudicial; en Argentina las diligencias preliminares son netamente judiciales.

\section{Concepto y naturaleza del proceso de desalojo}

El desalojo busca la \textbf{restitución de la tenencia} frente a quien carece de derecho a retenerla. Dentro del CPCCN es un proceso especial. Puede tramitar:

\begin{center}
\begin{tabularx}{\textwidth}{>{\raggedright\arraybackslash}p{4.2cm} Y}
\toprule
\textbf{Jurisdicción} & \textbf{Procesos de conocimiento disponibles} \\
\midrule
\textbf{Orden nacional} & Ordinario y sumarísimo (el sumario fue eliminado por la Ley 25.488) \\
\addlinespace
\textbf{Provincia de Buenos Aires} & Ordinario, sumario y sumarísimo \\
\bottomrule
\end{tabularx}
\end{center}

Es el juez quien decide la vía, y en la práctica la \textbf{mayoría tramita por sumarísimo} (debate más acotado). Pero llamarlo sumarísimo no garantiza celeridad real: la cátedra recuerda que los desalojos pueden durar años. En el sumarísimo, la contestación de demanda se rinde, por ejemplo, en 5 días.

\begin{note}
\textbf{Contexto histórico y comparado.} Un modelo más ágil es el \textbf{proceso monitorio} (otras provincias): ver Unidad 10. Córdoba tiene un desalojo abreviado con fuerte mirada sobre la parte vulnerable. Cada provincia tiene su código procesal; los hay muy actualizados (por ejemplo, el de San Juan).
\end{note}

\section{Enfoque sistémico y debido proceso}

\begin{itemize}
  \item \textbf{Principios} (creados por el constituyente: Constitución y tratados con jerarquía constitucional desde 1994) vs.\ \textbf{sistemas} (creados por el legislador para bajar el principio a la realidad). Si el sistema no respeta el principio, se plantean inconstitucionalidad e inconvencionalidad.
  \item \textbf{Debido proceso legal:} faceta \textbf{adjetiva} (notificación correcta, traslado, plazo de defensa) y \textbf{sustantiva} (respeto a la Constitución y tratados de derechos humanos).
  \item \textbf{Reglas de Brasilia} sobre acceso a la justicia de personas en condición de vulnerabilidad (Acordada 5/2009 CSJN): aplican a todo proceso. La vulnerabilidad no es una categoría fija; depende del contexto. Un jubilado que vive del alquiler de un pequeño departamento también puede ser vulnerable: hay que mirar \textbf{ambos lados de la puerta}.
\end{itemize}

\section{Acción, pretensión y demanda}

La \textbf{acción} es el derecho a acudir a la jurisdicción; la \textbf{pretensión} es lo que concretamente se reclama; la \textbf{demanda} es el acto que inicia el proceso. El \textbf{objeto} de la pretensión tiene tres elementos: \textbf{subjetivo} (quién reclama y a quién), \textbf{objetivo} (el tipo de sentencia que se pide) y la \textbf{causa} (hechos y derecho invocados). Normas aplicables: CCyC (fondo) y CPCCN (forma).

\section{Preclusión e inactividad procesal}

Rige la \textbf{preclusión}: cada etapa cumplida o vencida queda cerrada (contrasta con discovery/disclosure, donde se exhiben hechos y pruebas antes del juicio). Las etapas del proceso civil son seis y sucesivas: \textbf{introductoria, preliminar, probatoria, conclusional, recursiva y ejecutoria}. La inactividad de la parte que debe impulsar el proceso puede llevar a la \textbf{caducidad de instancia} (los apuntes solo la mencionan).

\section{Competencia, juez y recusación}

\begin{itemize}
  \item Primero se lee el contrato: ¿hay \textbf{foro prorrogado} o cláusula compromisoria?
  \item Si no, rigen los arts.\ 5 y ss.\ CPCCN: \textbf{juez del lugar del inmueble} o, subsidiariamente, el del domicilio del demandado.
  \item Puede haber \textbf{prórroga tácita} si el demandado contesta sin oponer incompetencia.
  \item La distribución de causas es de orden público: no se elige el juzgado; solo cabe \textbf{recusación con causa} (art.\ 17). La recusación sin causa está excluida (art.\ 14).
\end{itemize}

\begin{important}
\textbf{Fuero de atracción.} Si el contrato de locación tiene \textbf{destino comercial} y el locatario entra en concurso o quiebra, opera el fuero de atracción. Si el destino es \textbf{habitacional}, el contrato queda ajeno al concurso y no opera.
\end{important}

\section{Cuadro Cornell}

\begin{longtable}{Q R}
\toprule
\textbf{Preguntas / palabras clave} & \textbf{Notas / respuestas} \\
\midrule
\endfirsthead
\toprule
\textbf{Preguntas / palabras clave} & \textbf{Notas / respuestas} \\
\midrule
\endhead

\textbf{¿Para qué sirve el art.\ 323, inc.\ 1, y en cuánto caduca la diligencia preliminar?} &
Permite pedir que quien será demandado preste declaración jurada, por escrito, sobre hechos relativos a su personalidad, lo que sirve para identificar al ocupante. No se usa para averiguar el domicilio. La diligencia caduca a los 30 días desde su sustanciación si no se inicia el proceso principal. \\
\addlinespace[0.6em]

\textbf{¿En qué se diferencia la diligencia preliminar de la prueba anticipada?} &
La prueba anticipada es un sistema asegurativo para evitar la pérdida de una prueba (por ejemplo, un testigo próximo a fallecer) y es un instituto autónomo; las diligencias preliminares preparan el proceso de conocimiento y son de carácter restrictivo y residual. \\
\addlinespace[0.6em]

\textbf{¿Qué procesos de conocimiento hay en el orden nacional y en Buenos Aires?} &
Orden nacional: ordinario y sumarísimo (el sumario fue eliminado por la Ley 25.488). Provincia de Buenos Aires: ordinario, sumario y sumarísimo. La mayoría de los desalojos tramita por sumarísimo, aunque eso no garantiza celeridad real. \\
\addlinespace[0.6em]

\textbf{¿Qué diferencia hay entre principios y sistemas?} &
Los principios los crea el constituyente (Constitución y tratados con jerarquía constitucional desde 1994); los sistemas los crea el legislador para bajar el principio a la realidad. Si el sistema no respeta el principio, se plantean inconstitucionalidad e inconvencionalidad. \\
\addlinespace[0.6em]

\textbf{¿Qué abarca el debido proceso legal y qué dicen las Reglas de Brasilia?} &
Una faceta adjetiva (notificación correcta, traslado, plazo de defensa) y una sustantiva (respeto a la Constitución y a los tratados de derechos humanos). Las Reglas de Brasilia (Acordada 5/2009 CSJN) regulan el acceso a la justicia de personas vulnerables y aplican a todo proceso; la vulnerabilidad depende del contexto y hay que mirar ambos lados de la puerta. \\
\addlinespace[0.6em]

\textbf{¿Qué son acción, pretensión y demanda, y qué elementos tiene el objeto de la pretensión?} &
Acción: derecho a acudir a la jurisdicción. Pretensión: lo que concretamente se reclama. Demanda: acto que inicia el proceso. El objeto tiene un elemento subjetivo (quién reclama y a quién), uno objetivo (tipo de sentencia que se pide) y la causa (hechos y derecho invocados). \\
\addlinespace[0.6em]

\textbf{¿Cuáles son las etapas del proceso civil?} &
Seis y sucesivas: introductoria, preliminar, probatoria, conclusional, recursiva y ejecutoria. Rige la preclusión: cada etapa cumplida o vencida queda cerrada. \\
\addlinespace[0.6em]

\textbf{¿Cómo se determina la competencia y qué dice el art.\ 14 sobre la recusación?} &
Primero se lee el contrato (foro prorrogado o cláusula compromisoria); si no hay, rigen los arts.\ 5 y ss.\ CPCCN: juez del lugar del inmueble o, subsidiariamente, el del domicilio del demandado. Puede haber prórroga tácita si el demandado contesta sin oponer incompetencia. Solo cabe recusación con causa (art.\ 17); el art.\ 14 excluye la recusación sin causa. \\
\addlinespace[0.6em]

\textbf{¿Cuándo opera el fuero de atracción?} &
Si el contrato de locación tiene destino comercial y el locatario entra en concurso o quiebra. Si el destino es habitacional, el contrato queda ajeno al concurso y no opera. \\

\midrule
\multicolumn{2}{p{0.92\textwidth}}{%
\textbf{Resumen:} Diligencia preliminar: art.\ 323; inc.\ 1 sirve para identificar al ocupante; caduca a los \textbf{30 días}. No es prueba anticipada. Nacional: ordinario o sumarísimo; PBA: también sumario. La mayoría va por sumarísimo. Principios (constituyente) vs.\ sistemas (legislador). Competencia: contrato $\rightarrow$ arts.\ 5 y ss.\ (lugar del inmueble). Art.\ 14: sin recusación sin causa.} \\
\bottomrule
\end{longtable}

% ============================================================
\chapter[Partes, subinquilinos, fiadores y terceros]{Partes, subinquilinos, fiadores y terceros}
% ============================================================

\ideacentral{Hay que distinguir tres cosas: ser parte (capacidad), ser el titular de la relación (legitimación) y poder actuar en nombre de otro (personería). Y en el desalojo siempre hay que pensar en quiénes ocupan realmente el inmueble.}

\section{Capacidad, calidad de parte y legitimación}

\begin{center}
\begin{tabularx}{\textwidth}{>{\raggedright\arraybackslash}p{3.4cm} Y Y}
\toprule
\textbf{Concepto} & \textbf{Qué mide} & \textbf{Ejemplo} \\
\midrule
\textbf{Legitimación sustancial} & Ser \textbf{titular de la relación jurídica de fondo} (capacidad de derecho) & Un menor puede ser titular de un inmueble \\
\addlinespace
\textbf{Legitimación procesal / calidad de parte} & Aptitud para actuar por sí en el proceso (capacidad de ejercicio, art.\ 24 CCyC) & Un menor no puede demandar por sí; necesita representante (salvo emancipado) \\
\addlinespace
\textbf{Personería} & Representación con la que se actúa (poderes, tutela, curatela, representación societaria) & Poder mal otorgado: falta de personería \\
\bottomrule
\end{tabularx}
\end{center}

Puede tenerse una sin la otra (por ejemplo, una persona declarada incapaz por sentencia). Defensas: si falla la capacidad o representación, \textbf{falta de personería}; si quien reclama no es el sujeto habilitado, \textbf{falta de legitimación}. En la práctica a veces se las nombra con imprecisión.

\begin{example}{El caso del tío}
Un menor es titular de un inmueble alquilado y su tío, que lo administraba de hecho, demanda por falta de pago por derecho propio. \textbf{No tiene legitimación}: el legitimado es el locador (el menor, representado por su padre o tutor). El tío solo podría actuar como abogado patrocinante o con poder del representante legal.
\end{example}

\section{Legitimación activa y pasiva}

\begin{itemize}
  \item \textbf{Activa (no está regulada expresamente):} es amplia. Puede demandar el locador (aunque no sea dueño), el locatario que sublocó, el sublocador, el comodante, el heredero, el condómino, el usufructuario, el administrador, el representante y hasta el \textbf{poseedor} frente a quien tenga menos derecho.
  \item \textbf{Pasiva (art.\ 680 CPCCN):} puede demandarse a cualquiera que use u ocupe el inmueble y deba restituirlo: locatario, sublocatario, intruso, ocupantes de cualquier especie y sucesores.
  \item Un amigo del locatario que no ocupa el inmueble carece de legitimación pasiva.
\end{itemize}

\subsection*{Demandar a intrusos: diferencia de jurisprudencia}

\begin{center}
\begin{tabularx}{\textwidth}{>{\raggedright\arraybackslash}p{2.6cm} Y Y}
\toprule
 & \textbf{Provincia de Buenos Aires (departamentos judiciales)} & \textbf{Fuero civil nacional} \\
\midrule
\textbf{Qué se exige al actor} & Acreditar que tuvo en algún momento la \textbf{posesión efectiva}; el informe de dominio no basta (acredita título, no posesión) & Basta acompañar el título y acreditar el carácter invocado \\
\addlinespace
\textbf{Riesgo} & La demanda puede rechazarse si no se acredita la posesión pretérita & Menor exigencia \\
\bottomrule
\end{tabularx}
\end{center}

Contra intrusos hay también vía penal, pero la cátedra la considera poco práctica; la vía civil incluye el desalojo inmediato (Unidad 8). Las vías no se excluyen, aunque van por fueros distintos.

\section{Subinquilinos y ocupantes}

Se demanda siempre a \textbf{«subinquilinos y ocupantes»} en forma genérica, aunque no se los identifique por nombre, y se los notifica a todos. Exigir identificación nominal vulneraría la defensa en juicio (art.\ 18 CN) de personas no individualizadas. Aun si el contrato prohíbe sublocar, se notifica igual.

\section{Fiadores}

\begin{itemize}
  \item La fianza es una \textbf{garantía personal} (no sigue a un bien, a diferencia de la hipoteca, que es real). Responde por las deudas del locatario, incluidos los gastos pactados.
  \item \textbf{Art.\ 1225 CCyC:} la responsabilidad del fiador se limita al plazo contractual. Si la locación sigue de hecho sin renovación formal ni nueva fianza expresa, el fiador puede liberarse por el período posterior, \textbf{si su abogado lo opone a tiempo}. Según la clase 6, se extiende solo si el locador intima fehacientemente la restitución y notifica esa intimación al fiador.
  \item \textbf{¿Se lo notifica del desalojo?} Para la clase 5 no es un deber legal sino una buena práctica que algunos juzgados adoptan porque apura el acuerdo. En la audiencia de la clase 11 se lo trata como \textbf{parte}: la litis debe integrarse, y si no concurrió se le notifica el acuerdo para que preste conformidad u oposición.
  \item \textbf{Cautelares sobre bienes del fiador} (inhibición general, prohibición de contratar): no integran el desalojo, que busca la tenencia. Juegan como subsistema propio, típicamente en la ejecución de alquileres, y suelen exigir litis trabada.
\end{itemize}

% TODO: contenido incompleto o ambiguo en las notas originales
\begin{important}
\textbf{Inconsistencia a verificar.} Sobre el fiador, la clase 5 dice que notificarlo no es obligatorio y la clase 11 lo describe como parte cuya presencia integra la litis. Ambas son compatibles si se entiende «buena práctica» frente a «integración» en la audiencia, pero el examen oral puede pedir precisión.
\end{important}

\section{Terceros, litisconsorcio y desalojo administrativo}

\begin{itemize}
  \item En el desalojo la regla es el \textbf{litisconsorcio voluntario}; el necesario (demandar a todos, como en la escrituración contra condóminos) no se presenta.
  \item Cargas, deberes y obligaciones: \textbf{carga} = imperativo del propio interés (contestar); \textbf{deber} = su incumplimiento se sanciona; \textbf{obligación} = nace de una relación concreta (entregar la tenencia tras una resolución firme).
  \item \textbf{Desalojo administrativo:} solo excepcional, por decreto, en emergencias extremas (por ejemplo, riesgo de derrumbe), con control judicial posterior vía amparo. Fuera de esas causales corresponde el Poder Judicial.
\end{itemize}

\section{Representación: patrocinio y apoderado}

\begin{center}
\begin{tabularx}{\textwidth}{Y Y}
\toprule
\textbf{Patrocinio letrado} & \textbf{Apoderado} \\
\midrule
Requiere siempre la firma de la parte junto con la del letrado; sin ella el escrito se tiene por no presentado & Representa directamente a la parte. Art.\ 47 CPCCN exigía escritura pública; desde 2015, el art.\ 1017 CCyC admite cualquier forma. El poder se acompaña en la primera presentación \\
\bottomrule
\end{tabularx}
\end{center}

\section{Cuadro Cornell}

\begin{longtable}{Q R}
\toprule
\textbf{Preguntas / palabras clave} & \textbf{Notas / respuestas} \\
\midrule
\endfirsthead
\toprule
\textbf{Preguntas / palabras clave} & \textbf{Notas / respuestas} \\
\midrule
\endhead

\textbf{¿Qué diferencia hay entre legitimación sustancial, legitimación procesal y personería?} &
Legitimación sustancial: ser titular de la relación jurídica de fondo (capacidad de derecho). Legitimación procesal o calidad de parte: aptitud para actuar por sí en el proceso (capacidad de ejercicio, art.\ 24 CCyC). Personería: la representación con la que se actúa (poderes, tutela, curatela, representación societaria). \\
\addlinespace[0.6em]

\textbf{¿Qué diferencia hay entre falta de legitimación y falta de personería?} &
Si falla la capacidad o la representación, hay falta de personería; si quien reclama no es el sujeto habilitado, hay falta de legitimación. En la práctica a veces se las nombra con imprecisión. En el caso del tío, quien demanda por derecho propio sin ser locador carece de legitimación. \\
\addlinespace[0.6em]

\textbf{¿Quién tiene legitimación activa en el desalojo y quién pasiva?} &
La activa no está regulada expresamente y es amplia: locador (aunque no sea dueño), locatario que sublocó, sublocador, comodante, heredero, condómino, usufructuario, administrador, representante y hasta el poseedor frente a quien tenga menos derecho. La pasiva (art.\ 680 CPCCN) alcanza a cualquiera que use u ocupe el inmueble y deba restituirlo: locatario, sublocatario, intruso, ocupantes de cualquier especie y sucesores. \\
\addlinespace[0.6em]

\textbf{¿Qué se exige para demandar a intrusos en Buenos Aires y en el fuero civil nacional?} &
En los departamentos judiciales de Buenos Aires se exige acreditar que el actor tuvo en algún momento la posesión efectiva (el informe de dominio no basta); si no, la demanda puede rechazarse. En el fuero civil nacional basta acompañar el título y acreditar el carácter invocado. \\
\addlinespace[0.6em]

\textbf{¿Por qué se demanda a «subinquilinos y ocupantes» en forma genérica?} &
Exigir identificación nominal vulneraría la defensa en juicio (art.\ 18 CN) de personas no individualizadas. Se los notifica a todos, aun si el contrato prohíbe sublocar. \\
\addlinespace[0.6em]

\textbf{¿Hasta cuándo responde el fiador según el art.\ 1225?} &
Su responsabilidad se limita al plazo contractual. Si la locación sigue de hecho sin renovación formal ni nueva fianza expresa, puede liberarse por el período posterior si su abogado lo opone a tiempo. Según la clase 6, se extiende solo si el locador intima fehacientemente la restitución y notifica esa intimación al fiador. \\
\addlinespace[0.6em]

\textbf{¿Se notifica al fiador del desalojo?} &
Para la clase 5 es una buena práctica y no un deber legal. En la audiencia de la clase 11 se lo trata como parte: la litis debe integrarse y, si no concurrió, se le notifica el acuerdo para que preste conformidad u oposición. \\
\addlinespace[0.6em]

\textbf{¿Qué diferencia hay entre carga, deber y obligación?} &
Carga: imperativo del propio interés (por ejemplo, contestar). Deber: su incumplimiento se sanciona. Obligación: nace de una relación concreta (por ejemplo, entregar la tenencia tras una resolución firme). \\
\addlinespace[0.6em]

\textbf{¿Qué diferencia hay entre patrocinio letrado y apoderado?} &
El patrocinio requiere la firma de la parte junto con la del letrado; sin ella el escrito se tiene por no presentado. El apoderado representa directamente a la parte; desde 2015, el art.\ 1017 CCyC admite cualquier forma de poder, que se acompaña en la primera presentación. \\

\midrule
\multicolumn{2}{p{0.92\textwidth}}{%
\textbf{Resumen:} Legitimación sustancial = titular de la relación; procesal = capacidad de actuar; personería = representación. No hace falta ser dueño para demandar: basta ser locador. Pasiva (art.\ 680): locatario, sublocatario, intruso, ocupantes. Subinquilinos y ocupantes: fórmula genérica por defensa en juicio. Fiador (art.\ 1225): responde hasta el plazo contractual; ojo con la prórroga de hecho.} \\
\bottomrule
\end{longtable}

